\section{Conclusão: A Nova Fronteira do Ensino Odontológico}
\label{sec:educacao-conclusao}

Os gêmeos digitais delineiam o limiar de um novo período evolutivo na educação odontológica, promovendo uma desconstrução sistemática das limitações espaciais, temporais e éticas inerentes ao modelo pedagógico centrado em extrações prévias e preceptorado presencial de observação e imitação. A convergência entre simulação háptica, modelagem biomecânica preditiva e ecossistemas \textit{web-based} provê um terreno infinitamente reconfigurável para a lapidação do erro iterativo e do aprendizado seguro sem exposição iatrogênica do paciente~\cite{techsoc2025precision}.

Essas plataformas, subsidiadas pelas arquiteturas multiagentes como as desenvolvidas no ecossistema OpenCode, assumem a responsabilidade de catalisadores da transição da arte empírica artesanal para a ciência de dados de alta precisão. A orquestração das evidências automatizadas, os repositórios ômicos de casos fenotipicamente curados e os algoritmos de tutoria neural garantem uma curva de assimilação técnica mais acentuada e uma cognição clínica infinitamente mais plural.

Os entraves de validação científica de longo prazo, da onerosidade de infraestrutura e da relutância metodológica representam atritos expectáveis frente a disrupções desse porte. Não obstante, o vetor direcional é inescapável. As instituições de ensino que se omitirem na integração dessas biotecnologias virtuais correm o risco sumário de obsolescência epistêmica em menos de uma década. Em última análise, a odontologia do futuro — sustentada por diagnósticos algorítmicos e manufaturas robóticas — exige um processo civilizatório educacional correlato: um ensino inerentemente digital, distribuído, imersivo e centrado na precisão cibernética.
